% Bogurodzica: performance edition
% Build:  gregorio bogurodzica-performance.gabc
%         lualatex bogurodzica-performance.tex   (twice; needs GregorioTeX 6 and Junicode)
\documentclass[12pt,a4paper]{article}
\usepackage[margin=2cm,top=1.8cm,bottom=2cm]{geometry}
\usepackage{fontspec}
\setmainfont{Junicode}[Path=/usr/share/fonts/opentype/junicode/,Extension=.otf,UprightFont=*-Regular,ItalicFont=*-Italic,BoldFont=*-SmBold,BoldItalicFont=*-SmBoldItalic,Numbers=OldStyle]
\usepackage{polyglossia}
\setdefaultlanguage[variant=british]{english}
\setotherlanguage{polish}
\usepackage{microtype}
\usepackage{gregoriotex}
\usepackage{array,booktabs}
\usepackage{enumitem}
\usepackage{xcolor}
\definecolor{rubric}{RGB}{160,30,30}
\pagestyle{empty}

\setlength{\parindent}{0pt}
\setlength{\parskip}{0.5em}

\gresetinitiallines{1}
\grechangestyle{initial}{\fontsize{48}{48}\selectfont}
\grechangestaffsize{22}
\grechangedim{baselineskip}{70pt plus 6pt minus 6pt}{scalable}

\newcommand{\pl}[1]{\foreignlanguage{polish}{#1}}
\newcommand{\ipa}[1]{{\small\color{black!70}[#1]}}

\begin{document}

\begin{center}
{\LARGE\textsc{Bogurodzica}}\\[0.15em]
{\small Polish, 13th--14th century \quad·\quad Kraków, Biblioteka Jagiellońska, MS 1619}
\end{center}
\vspace{0.6em}

\input{bogurodzica-performance.gtex}

\newpage
\section*{\large Text, pronunciation and translation}

{\small Pronunciation follows modern Polish, the usual practice for this
song. \textit{y} is a vowel between English \textit{i} and \textit{u}; \textit{ł}
is English \textit{w}; \textit{ż}, \textit{sz} and \textit{cz} are hard
\textit{zh}, \textit{sh} and \textit{ch}; \textit{ś}, \textit{ć} and \textit{dzi}
are softer, closer to \textit{sh}, \textit{ch} and \textit{j}; \textit{ę} and
\textit{ą} are nasal. Sing \textit{Kyrie eleison} in seven syllables, as
underlaid.\par}

\renewcommand{\arraystretch}{1.0}\small
\begin{tabular}{@{}>{\raggedright\arraybackslash}p{0.62\textwidth}>{\raggedright\arraybackslash\itshape}p{0.34\textwidth}@{}}
\toprule
\multicolumn{2}{@{}l}{\textbf{1}}\\
\pl{Bogurodzica dziewica, Bogiem sławiena Maryja,}\newline\ipa{bɔɡurɔdʑitsa dʑɛvitsa bɔɡʲɛm swavʲɛna marɨja}
 & Mother of God, Virgin glorified by God, Mary!\\[0.15em]
\pl{u twego syna Gospodzina matko zwolena, Maryja!}\newline\ipa{u tfɛɡɔ sɨna ɡɔspɔdʑina matkɔ zvɔlɛna marɨja}
 & Chosen Mother, Mary, from thy Son the Lord\\[0.15em]
\pl{Zyszczy nam, spuści nam.}\newline\ipa{zɨʂtʂɨ nam spuɕtɕi nam}
 & win for us, send down to us.\\[0.15em]
\pl{Kyrie eleison.}\newline\ipa{kiriɛ ɛlɛisɔn}
 & Lord, have mercy.\\
\midrule
\multicolumn{2}{@{}l}{\textbf{2}}\\
\pl{Twego dziela Krzciciela, Bożyce,}\newline\ipa{tfɛɡɔ dʑɛla kʂtɕitɕɛla bɔʐɨtsɛ}
 & For thy Baptist's sake, Son of God,\\[0.15em]
\pl{usłysz głosy, napełń myśli człowiecze.}\newline\ipa{uswɨʂ ɡwɔsɨ napɛwɲ mɨɕli tʂwɔvʲɛtʂɛ}
 & hear our voices, fulfil the thoughts of men.\\[0.15em]
\pl{Słysz modlitwę, jąż nosimy,}\newline\ipa{swɨʂ mɔdlitfɛ̃ jɔ̃ʂ nɔɕimɨ}
 & Hear the prayer we bring,\\[0.15em]
\pl{a dać raczy, jegoż prosimy:}\newline\ipa{a datɕ ratʂɨ jɛɡɔʂ prɔɕimɨ}
 & and deign to grant what we ask:\\[0.15em]
\pl{a na świecie zbożny pobyt,}\newline\ipa{a na ɕfʲɛtɕɛ zbɔʐnɨ pɔbɨt}
 & on earth a godly life,\\[0.15em]
\pl{po żywocie rajski przebyt.}\newline\ipa{pɔ ʐɨvɔtɕɛ rajski pʂɛbɨt}
 & after life a dwelling in paradise.\\[0.15em]
\pl{Kyrie eleison.} & Lord, have mercy.\\
\bottomrule
\end{tabular}

\section*{\large Performance notes}
{\small\begin{itemize}[leftmargin=1.2em,itemsep=0.15em]
\item \textbf{Pitch.} The written final is D. The range is a ninth, C to the
  D above, so the score suits most voices as written; transpose freely.
\item \textbf{Rhythm.} The manuscript gives pitches only. Sing in the rhythm
  of the words, with notes of roughly equal length. Sing the notes of a group
  on one syllable lightly and evenly, and lengthen the last note before each
  bar line. The pace is unhurried.
\item \textbf{Breathing.} Breathe at the full bars. The half and quarter bars
  mark the rhymes inside each line: a short lift, or a catch breath if needed.
  The double bars end each stanza.
\item \textbf{The B on \pl{spuści}} is natural. It is the only B in the song,
  at the top of the melody. Some 16th-century copies flatten B elsewhere; this
  version has none.
\item \textbf{The leap at \pl{U twego}.} The line jumps an octave from the
  low D of \pl{Maryja} to the high D. Later sources smooth it out, but it is
  in the oldest copy. Take it cleanly.
\item \textbf{Forces.} The song was sung in church and, by Długosz's account,
  by Polish knights before battle, Grunwald 1410 among them. A single voice,
  a small schola or a full choir in unison all suit it. One option is a
  cantor for each stanza with everyone joining at \pl{Kyrie eleison}.
\item \textbf{Both stanzas.} Stanza 1 is addressed to the Virgin, stanza 2 to
  Christ. They share most of their music.
\end{itemize}}

\vfill
{\footnotesize Melody and text from the earliest notated copy, Kraków,
Biblioteka Jagiellońska, MS 1619, written shortly after 1400. Text in
normalised Old Polish after J. Woronczak (1962). Bar lines are editorial. For
sources and variants see the critical edition.\par}

\end{document}
